\par\medskip
\hypertarget{native-075-article057}{}\pdfbookmark[2]{57.}{article-057}
\begin{center}57.\end{center}
Corporis itaque coelestis distantia a quovis plano per Solem transeunte reduci poterit ad formam $kr\sin(v + K)$, designante $v$ anomaliam veram, eritque $k$ sinus inclinationis orbitae ad hoc planum, $K$ distantia perihelii a nodo ascendente orbitae in eodem plano. Quatenus situs plani orbitae, lineaeque apsidum in eo, nec non situs plani ad quod distantiae referuntur pro constantibus haberi possunt, etiam $k$ et $K$ constantes erunt. Frequentius tamen illa methodus in tali casu in usum vocabitur, ubi tertia saltem suppositio non permittitur, etiamsi perturbationes negligantur, quae primam atque secundam semper aliquatenus afficiunt. Illud evenit, quoties distantiae referuntur ad
\SourceLeaf{076}aequatorem, sive ad planum aequatorem sub angulo recto in
rectascensione data secans: quum enim situs aequatoris propter
praecessionem aequinoctiorum insuperque propter nutationem (siquidem
de vero non de medio situ sermo fuerit) mutabilis sit, in hoc casu
etiam $k$ et $K$ mutationibus, lentis utique, obnoxiae erunt.
Computus harum mutationum per formulas differentiales absque
difficultate eruendas absolvi potest: hic vero brevitatis caussa
sufficiat, variationes differentiales ipsarum $i'$, $\AscendingNode'$,
$\Delta$ apposuisse, quatenus a variationibus ipsarum $\AscendingNode - n$
atque $\epsilon$ pendent.
\[
\mathrm{d}i' = \sin\epsilon \sin\AscendingNode'\, \mathrm{d}(\AscendingNode - n) - \cos\AscendingNode'\, \mathrm{d}\epsilon
\]
\[
\mathrm{d}\AscendingNode' = \frac{\sin i \cos\Delta}{\sin i'}\, \mathrm{d}(\AscendingNode - n) + \frac{\sin\AscendingNode'}{\operatorname{tang} i'}\, \mathrm{d}\epsilon
\]
\[
\mathrm{d}\Delta = \frac{\sin\epsilon \cos\AscendingNode'}{\sin i'}\, \mathrm{d}(\AscendingNode - n) + \frac{\sin\AscendingNode'}{\sin i'}\, \mathrm{d}\epsilon.
\]
Ceterum quoties id tantum agitur, ut plures corporis coelestis loci
respectu talium planorum mutabilium calculentur, qui temporis
intervallum mediocre complectuntur (e. g. unum annum), plerumque
commodissimum erit, quantitates $a$, $A$, $b$, $B$, $c$, $C$ pro
duabus epochis intra quas illi cadunt reipsa calculare, ipsarumque
mutationes pro singulis temporibus propositis ex illis per simplicem
interpolationem eruere.
